% 两个点的四种标记及区间类的实现方式。
\begin{tikzpicture}[x=1mm,y=1mm,font=\footnotesize]
  \foreach \sx/\sy/\a/\b/\left/\right in {
    0/38/0/0/4/6,
    53/38/1/0/8/14,
    0/0/0/1/26/32,
    53/0/1/1/8/32}{
    \begin{scope}[xshift=\sx mm,yshift=\sy mm]
      \fill[BookPaper] (0,0) rectangle (42,27);
      \draw[BookMuted,thin] (5,8) -- (37,8);
      \ifnum\a=1
        \fill[FigureInput] (14,8) circle (2);
      \else
        \draw[FigureLoss,thick] (14,8) circle (2);
      \fi
      \ifnum\b=1
        \fill[FigureInput] (28,8) circle (2);
      \else
        \draw[FigureLoss,thick] (28,8) circle (2);
      \fi
      \draw[FigureModel,very thick] (\left,14) -- (\right,14);
      \draw[FigureModel,very thick] (\left,12.5) -- (\left,15.5);
      \draw[FigureModel,very thick] (\right,12.5) -- (\right,15.5);
      \node[font=\kaishu\scriptsize] at (21,22) {标记$(\a,\b)$};
    \end{scope}
  }
  \node[font=\kaishu\scriptsize,text=BookMuted,align=center] at (47,75)
    {实心点表示标签1，空心点表示标签0；蓝色线段是选取的区间。};
\end{tikzpicture}
